\documentclass[11pt, a4paper]{article}

% --- Paquetes Fundamentales ---
\usepackage[utf8]{inputenc}
\usepackage[T1]{fontenc}
\usepackage[spanish,es-tabla]{babel}
\usepackage[margin=2.5cm, headheight=15pt]{geometry}
\usepackage{graphicx}
\usepackage{fancyhdr}
\usepackage{xcolor}
\usepackage{microtype}
\usepackage[colorlinks=true, linkcolor=blue, urlcolor=blue]{hyperref}

% --- Matemáticas y Electrónica ---
\usepackage{amsmath, amssymb}
\usepackage{siunitx}
\usepackage[american, RPvoltages]{circuitikz}

% --- Elementos de Diseño Pedagógico y Tablas ---
\usepackage{tcolorbox}
\usepackage{tabularx}
\usepackage{booktabs}
\usepackage{colortbl}

% --- Configuración de Colores y siunitx ---
\definecolor{ucbblue}{RGB}{0, 51, 102}
\sisetup{
    output-decimal-marker = {,},
    per-mode = symbol,
    inter-unit-product = \ensuremath{{}\cdot{}}
}

% --- Macros y Estilos ---
\tcbset{
    caja_conclusion/.style={colback=blue!3, colframe=ucbblue, fonttitle=\bfseries, boxrule=0.8pt, arc=3pt}
}

% --- Estilo de Página ---
\pagestyle{fancy}
\fancyhf{}
\lhead{\textcolor{ucbblue}{\textbf{IMT-121 | Resolución Analítica}}}
\rhead{\textbf{Análisis de Transferencia (.tf)}}
\cfoot{Página \thepage}
\renewcommand{\headrulewidth}{0.4pt}
\renewcommand{\headrule}{\hbox to\headwidth{\color{ucbblue}\leaders\hrule height \headrulewidth\hfill}}

\begin{document}

\begin{center}
    {\LARGE \textbf{Resolución Analítica: Parámetros de Transferencia en CC}} \\[0.3cm]
    {\large \textit{Validación Matemática del Análisis \texttt{.tf} en LTspice}}
    \vspace{0.4cm}
\end{center}
\hrule
\vspace{0.6cm}

Para determinar los valores exactos que arroja LTspice en el análisis \texttt{.tf} (Transfer Function) para este circuito, evaluamos analíticamente paso a paso cada parámetro:

\begin{center}
    \begin{circuitikz}[scale=1, transform shape]
        % Fuente de entrada y R1
        \draw (0,0) to[V, v=$V_1 {=} \qty{10}{\volt}$] (0,3)
              to[R, l=$R_1 {=} \qty{50}{\ohm}$] (4,3) node[above=0.1cm]{\texttt{N002}};
        % VCVS invertida para que el (+) esté en el N003
        \draw (4,3) to[cV, l_=$E_1$, v=$3V_x$, invert] (8,3) node[above=0.1cm]{\texttt{N003}};
        % Ramas a tierra
        \draw (4,3) to[R, l=$R_2 {=} \qty{100}{\ohm}$, v=$V_x$] (4,0);
        \draw (8,3) to[R, l=$R_3 {=} \qty{200}{\ohm}$] (8,0);
        % GND
        \draw (0,0) -- (8,0);
        \draw (4,0) node[ground]{};
    \end{circuitikz}
\end{center}

\section*{1. Cálculo de la Función de Transferencia (\texttt{Transfer\_function})}
En LTspice, la fuente dependiente \texttt{e} colocada entre el Nodo 2 y el Nodo 3 con el terminal positivo hacia la salida define:
\begin{equation*}
    V(\text{N003}) - V(\text{N002}) = 3 \cdot V_x
\end{equation*}
Dado que la variable de control es la tensión del Nodo 2 ($V_x = V(\text{N002})$):
\begin{equation*}
    V(\text{N003}) - V(\text{N002}) = 3 \cdot V(\text{N002}) \implies V(\text{N003}) = 4 \cdot V(\text{N002})
\end{equation*}

Aplicando la Ley de Corrientes de Kirchhoff en el supernodo formado por los Nodos 2 y 3:
\begin{equation*}
    \frac{V(\text{N002}) - V_1}{R_1} + \frac{V(\text{N002})}{R_2} + \frac{V(\text{N003})}{R_3} = 0
\end{equation*}

Sustituyendo los valores ($R_1 = \qty{50}{\ohm}$, $R_2 = \qty{100}{\ohm}$, $R_3 = \qty{200}{\ohm}$ y $V(\text{N003}) = 4 V(\text{N002})$):
\begin{align*}
    \frac{V(\text{N002}) - V_1}{50} + \frac{V(\text{N002})}{100} + \frac{4 V(\text{N002})}{200} &= 0 \\
    \frac{V(\text{N002}) - V_1}{50} + \frac{3 V(\text{N002})}{100} &= 0
\end{align*}

Multiplicando por $100$:
\begin{align*}
    2(V(\text{N002}) - V_1) + 3 V(\text{N002}) &= 0 \\
    5 V(\text{N002}) &= 2 V_1 \implies V(\text{N002}) = 0.4 \cdot V_1
\end{align*}

Con $V_1 = \qty{10}{\volt}$:
\begin{align*}
    V(\text{N002}) &= \qty{4.0}{\volt} \\
    V(\text{N003}) &= 4 \cdot (\qty{4.0}{\volt}) = \qty{16.0}{\volt}
\end{align*}

\begin{equation*}
    \text{Transfer\_function} = \frac{V(\text{N003})}{V_1} = \frac{\qty{16.0}{\volt}}{\qty{10.0}{\volt}} = \mathbf{1.600}
\end{equation*}

\section*{2. Cálculo de la Impedancia de Entrada (\texttt{v1\#Input\_impedance})}
La corriente entregada por la fuente independiente $V_1$ es:
\begin{equation*}
    I_{\text{in}} = \frac{V_1 - V(\text{N002})}{R_1} = \frac{\qty{10}{\volt} - \qty{4}{\volt}}{\qty{50}{\ohm}} = \frac{\qty{6}{\volt}}{\qty{50}{\ohm}} = \qty{0.12}{\ampere} = \qty{120}{\milli\ampere}
\end{equation*}

La impedancia de entrada vista por la fuente $V_1$ corresponde a:
\begin{equation*}
    R_{\text{in}} = \frac{V_1}{I_{\text{in}}} = \frac{\qty{10}{\volt}}{\qty{0.12}{\ampere}} = \frac{250}{3}\,\Omega = \mathbf{\qty{83.333}{\ohm}}
\end{equation*}

\section*{3. Cálculo de la Impedancia de Salida (\texttt{output\_impedance\_at\_v(n003)})}
Para hallar la resistencia de salida (Thévenin en el Nodo 3), se apaga la fuente independiente ($V_1 = \qty{0}{\volt}$) y se inyecta una corriente de prueba $I_{\text{test}}$ en el Nodo 3:

Al apagar $V_1$, la resistencia equivalente desde el Nodo 2 hacia tierra es:
\begin{equation*}
    R_p = R_1 \parallel R_2 = \frac{50 \cdot 100}{50 + 100} = \frac{100}{3}\,\Omega \approx \qty{33.333}{\ohm}
\end{equation*}

La corriente $I_{\text{test}}$ se divide entre la resistencia $R_3$ y la rama que regresa por la fuente controlada $E_1$:
\begin{equation*}
    I_{\text{test}} = \frac{V_{\text{test}}}{R_3} + I_{E1}
\end{equation*}

Como $V_{\text{test}} = 4 V(\text{N002}) \implies V(\text{N002}) = \frac{V_{\text{test}}}{4}$, la corriente que fluye por la rama de $E_1$ hacia el paralelo $R_p$ es:
\begin{equation*}
    I_{E1} = \frac{V(\text{N002})}{R_p} = \frac{V_{\text{test}} / 4}{100 / 3} = \frac{3}{400} V_{\text{test}} = 0.0075 \cdot V_{\text{test}}
\end{equation*}

La resistencia de dicha rama es $R_{\text{rama}} = \frac{V_{\text{test}}}{I_{E1}} = \frac{400}{3}\,\Omega \approx \qty{133.333}{\ohm}$.

La resistencia total de salida es el paralelo entre $R_3$ ($\qty{200}{\ohm}$) y $R_{\text{rama}}$:
\begin{equation*}
    R_{\text{out}} = 200 \parallel \frac{400}{3} = \frac{200 \cdot \frac{400}{3}}{200 + \frac{400}{3}} = \frac{\frac{80000}{3}}{\frac{1000}{3}} = \mathbf{\qty{80.000}{\ohm}}
\end{equation*}

\vspace{0.4cm}
\section*{Valores de Salida en la Ventana de LTspice}
Al ejecutar la directiva \texttt{.tf V(N003) V1}, el cuadro de resultados de LTspice muestra exactamente:

\vspace{0.2cm}
\begin{table}[h!]
    \centering
    \renewcommand{\arraystretch}{1.5}
    \begin{tabularx}{\textwidth}{@{}l c X@{}}
        \toprule
        \rowcolor{ucbblue!10}
        \textbf{Parámetro en LTspice} & \textbf{Valor Reportado} & \textbf{Interpretación Física} \\
        \midrule
        \texttt{Transfer\_function} & \texttt{1.600000e+00} & Ganancia de tensión $\frac{V_3}{V_1} = \qty{1.6}{\volt\per\volt}$ \\
        \texttt{output\_impedance\_at\_v(n003)} & \texttt{8.000000e+01} & Resistencia de Thévenin en la salida = $\qty{80}{\ohm}$ \\
        \texttt{v1\#Input\_impedance} & \texttt{8.333333e+01} & Resistencia de entrada vista por la fuente = $\qty{83.33}{\ohm}$ \\
        \bottomrule
    \end{tabularx}
\end{table}

\vspace{0.4cm}


\end{document}